\documentclass{article}

\input{../definitions.tex}

\usepackage{mathpartir}
\usepackage{tabto}

\newdef{problem}
\newdef{exercise}
% \usepackage[a4paper,margin=1cm,landscape]{geometry}
\newpf{solution}{}
\usepackage[a4paper,margin=2cm]{geometry}
%\NewDocumentEnvironment{solution}{+b}{}{}

\ExplSyntaxOn
\NewDocumentEnvironment{problemdef}{o +b}
{\tl_gset:Nn \g_tmpa_tl {#2}}
{
  \begin{problem}
    #1
    \tl_use:N \g_tmpa_tl
  \end{problem}
  \begin{solution}
    \DeclareDocumentCommand{\emptybox}{m}{\boxed{##1}}
    \tl_use:N \g_tmpa_tl
  \end{solution}
}
\ExplSyntaxOff

\newcommand{\seq}[3]{#1 ⊢ #2 : #3}
\def\emp{\bullet}
\newcommand{\lam}[3]{λ #1 : #2.\: #3}
\renewcommand{\pair}[2]{\p{#1, #2}}
\newcommand{\pr}[2]{\textsf{pr}_{#1} {#2}}
\newcommand{\match}[5]{\textsf{match}\:{#1}\:| \inl {#2} ↦ {#3} | \inr {#4} ↦ {#5}}
\newcommand{\emptybox}[1]{\boxed{\phantom{#1}}}

% \def\trees{
% \newcommand{\nullaryrule}[2]{
%   \ExpandArgs{c}\DeclareDocumentCommand{#1}{m}{\inferrule*[right=\rlap{\(#2\)}] { } {##1}}
% }
% \newcommand{\unaryrule}[2]{
%   \ExpandArgs{c}\DeclareDocumentCommand{#1}{mm}{\inferrule*[right=\rlap{\(#2\)}] {##2} {##1}}
% }
% \newcommand{\binaryrule}[2]{
%   \ExpandArgs{c}\DeclareDocumentCommand{#1}{mmm}{\inferrule*[right=\rlap{\(#2\)}] {##2 \\ ##3} {##1}}
% }
% \newcommand{\ternaryrule}[2]{
%   \ExpandArgs{c}\DeclareDocumentCommand{#1}{mmmm}{\inferrule*[right=\rlap{\(#2\)}] {##2 \\ ##3 \\ ##4} {##1}}
% }
% \nullaryrule{ax}{\textsc{Ax}}
% \binaryrule{prodI}{×\textsc{I}}
% \unaryrule{prodEL}{×\textsc{E}₁}
% \unaryrule{prodER}{×\textsc{E}₂}
% \unaryrule{funI}{→\textsc{I}}
% \binaryrule{funE}{→\textsc{E}}
% \unaryrule{sumIL}{+\textsc{I}{}₁}
% \unaryrule{sumIR}{+\textsc{I}{}₂}
% \ternaryrule{sumE}{+\textsc{E}}
% }
\newcommand{\linear}{
\newcommand{\nullaryrule}[2]{
  \ExpandArgs{c}\DeclareDocumentCommand{##1}{mO{}}{\item \(####1\) \tab \(##2\) ####2}
}
\newcommand{\unaryrule}[2]{
  \ExpandArgs{c}\DeclareDocumentCommand{##1}{mmO{}}{{####2} {\item \(####1\) \tab \(##2\) ####3}}
}
\newcommand{\binaryrule}[2]{
  \ExpandArgs{c}\DeclareDocumentCommand{##1}{mmmO{}}{{####2} {####3} {\item \(####1\) \tab \(##2\) ####4}}
}
\nullaryrule{ass}{\text{assumption}}
\binaryrule{prodI}{∧\textsc{I}}
\unaryrule{prodEL}{∧\textsc{E}₁}
\unaryrule{prodER}{∧\textsc{E}₂}
\DeclareDocumentCommand{\funI}{smmmO{}}{
  {{\ass {##2}} {##4}}
  {\item \(\IfBooleanTF{##1}{\p{##2}}{##2} → ##3\) \tab \(→\textsc{I}\) ##5}}
\binaryrule{funE}{→\textsc{E}}
\unaryrule{sumIL}{∨\textsc{I}₁}
\unaryrule{sumIR}{∨\textsc{I}₂}
\DeclareDocumentCommand{\sumE}{mmmmmmO{}}{
  {##2}
  {{\ass {##3}} {##4}}
  {{\ass {##5}} {##6}}
  {\item \(##1\) \tab \(∨\textsc{E}\) ##7}}
\nullaryrule{botE}{⊥\textsc{E}}
\unaryrule{forE}{∀\textsc{E}}
\DeclareDocumentCommand{\forI}{mmmO{}}{
  {{\ass {##1}} {##3}}
  {\item \(\for{##1}{##2}\) \tab \(∀\textsc{I}\) ##4}}
\DeclareDocumentCommand{\existE}{mmmmmO{}}{
  {##2}
  {{\ass {##3}} {\ass {##4} {##5}}}
  {\item \(##1\) \tab \(∃\textsc{E}\) ##6}}
}

\begin{document}

\section{Linear-style natural deduction for predicate logic}

\subsection{Problems}

\begin{problemdef}
  [Show \(P → Q → P∧Q\) is valid by filling in the below proof.\\]
  \parbox[t][][t]{\textwidth}{
    \begin{enumerate}
    \item \(P\) \tab assumption
    \item \emptybox{Q} \tab \emptybox{\text{assumption}}
    \item \(P ∧ Q\) \tab \(∧\textsc{I}\) (1) (\emptybox{2})
    \item \(Q → P∧Q\) \tab \(→\textsc{I}\) (2--3)
    \item \(\emptybox{P → Q → P∧Q}\) \tab \(→\textsc{I}\) (\emptybox{\text{1--4}})
    \end{enumerate}
  }
\end{problemdef}


\begin{problem}
  Show \(P∧Q → Q∧P\) is valid.
\end{problem}
\begin{solution}
  \linear
  \parbox[t][][t]{\textwidth}{
    \begin{enumerate}
      \funI
      {P∧Q} {Q∧P}
      {\prodI
        {Q∧P}
        {\prodER {Q} {} [(1)]}
        {\prodEL {P} {} [(1)]}
        [(2) (3)]
      }
      [(1--4)]
    \end{enumerate}
  }
\end{solution}

\begin{problem}
  Show \(\p{P → Q → R} → Q → P → R\) is valid.
\end{problem}
\begin{solution}
  \linear
  \parbox[t][][t]{\textwidth}{
  \begin{enumerate}
    \funI*
    {P → Q → R} {Q → P → R}
    {\funI
      {Q} {P → R}
      {\funI
        {P} {R}
        {\funE
          {R}
          {\funE {Q → R} {} {} [(1) (3)]}
          {}
          [(4) (2)]
        }
        [(3--5)]
      }
      [(2--6)]
    }
    [(1--7)]
  \end{enumerate}
  }
\end{solution}


\begin{problem}
  Show \(P∨Q → Q∨P\) is valid.
\end{problem}
\begin{solution}
  \linear
  \parbox[t][][t]{\textwidth}{
    \begin{enumerate}
      \funI
      {P∨Q} {Q∨P}
      {\sumE
        {Q ∨ P}
        {}
        {P}
        {\sumIR
          {Q∨P}
          {}
          [(2)]
        }
        {Q}
        {\sumIL
          {Q∨P}
          {}
          [(4)]
        }
        [(1) (2--3) (4--5)]
      }
      [(1) (2--6)]
    \end{enumerate}
  }
\end{solution}


\begin{problem}
  Show \(P∧(Q∨R) → P∧Q ∨ P∧R\) is valid.
\end{problem}
\begin{solution}
  \linear
  \parbox[t][][t]{\textwidth}{
    \begin{enumerate}
      \funI
      {P∧(Q∨R)} {P∧Q ∨ P∧R}
      {\prodEL {P} {} [(1)]
        \sumE
        {P∧Q ∨ P∧R}
        {\prodEL {Q∨R} {} [(1)]}
        {Q}
        {\sumIL
          {P∧Q ∨ P∧R}
          {\prodI {P∧Q} {} {} [(2) (4)]}
          [(5)]
        }
        {R}
        {\sumIR
          {P∧Q ∨ P∧R}
          {\prodI {P∧R} {} {} [(2) (7)]}
          [(8)]
        }
        [(3) (4--6) (7--9)]
      }
      [(1--10)]
    \end{enumerate}
  }
\end{solution}

\begin{problem}
  Show \(P ∨ ⊥ → P\).
\end{problem}
\begin{solution}
  \linear
  \parbox[t][][t]{\textwidth}{
    \begin{enumerate}
      \funI
      {P ∨ ⊥} {P}
      {\sumE
        {P}
        {}
        {P}
        {}
        {⊥}
        {\botE P [(P)]}
        [(1) (2--2) (3--4)]
      }
      [(1--5)]
    \end{enumerate}
  }
\end{solution}

\begin{problem}
  Show \(P∨\for{x : A}{Q x} → \for{y : A}{P ∨ Qy}\).
\end{problem}
\begin{solution}
  \linear
  \parbox[t][][t]{\textwidth}{
    \begin{enumerate}
      \funI
      {P∨\for{x : A}{Q x}} {\for{x : A}{P ∨ Qx}}
      {\forI
        {y : A} {P ∨ Qy}
        {\sumE
          {P ∨ Qy}
          {}
          {P}
          {\sumIL {P ∨ Qy} {} [(3)]}
          {\for{x : A}{Q x}}
          {\sumIR
            {P ∨ Qy}
            {\forE {Q y} {} [(5) (2)]}
            [(6)]
          }
          [(1) (3--4) (5--7)]
        }
        [(2--8)]
      }
      [(1--9)]
    \end{enumerate}
  }
\end{solution}

\begin{problem}
  Show \(\for{x : A}{P x → Q} → \p{\exist{y : A}{Py}} → Q\).
\end{problem}
\begin{solution}
  \linear
  \parbox[t][][t]{\textwidth}{
    \begin{enumerate}
      \funI
      {\for{x : A}{P x → Q}} {\p{\exist{y : A}{Py}} → Q}
      {\funI
        {\exist{y : A}{Py}} {Q}
        {\existE
          {Q}
          {}
          {y : A}
          {Py}
          {\funE
            {Q}
            {}
            {\forE
              {Py → Q}
              {}
              [(1) (3)]
            }
            [(5) (4)]
          }
          [(2) (3,4--6)]
        }
        [(2--7)]
      }
      [(1--8)]
    \end{enumerate}
  }
\end{solution}

\begin{exercise}
  Show the converse to the above, aka \(\p{\p{\exist{y : A}{Py}} → Q} → \for{x : A}{P x → Q}\).
\end{exercise}
\begin{exercise}
  Compare the proofs here with the proofs (of the ``same'' problems) from last
  week. For the problems, which we did not give a tree-style proof last week,
  try it yourself. You can use the linear-style proofs as a blueprint.
\end{exercise}

\end{document}


%%% Local Variables:
%%% mode: latex
%%% TeX-master: t
%%% End:
